\chapter{Hasil Lengkap Studi Pendahuluan}\label{lam:hasil}

Tabel~\ref{tab:per-kelas} menyajikan F1 per kelas hasil replikasi kelompok (K), yaitu
rata-rata tiga seed, berdampingan dengan Tabel~3 \textcite{audebert2020} (P).

Selisih terbesar ada pada kelas Resume untuk model fusion: 0,803 berbanding 0,96. Selisih
inilah yang menjadi petunjuk pertama bahwa cabang teks tidak terpakai, karena model teks
saja mencapai 0,969 pada kelas yang sama.

\begin{table}[H]
  \centering
  \caption{F1 per kelas: replikasi kelompok (K) dan paper (P)}
  \label{tab:per-kelas}
  \fontsize{10}{12}\selectfont
  \begin{tabular}{|l|c|c|c|c|c|c|c|c|}
    \hline
    & \multicolumn{2}{c|}{\textbf{Teks}} & \multicolumn{2}{c|}{\textbf{Citra}} &
      \multicolumn{2}{c|}{\textbf{Fusion}} & \multicolumn{2}{c|}{\textbf{Oracle}} \\ \hline
    \textbf{Kelas} & K & P & K & P & K & P & K & P \\ \hline
    Advertisement & 0,537 & 0,60 & 0,902 & 0,94 & 0,894 & 0,93 & 0,930 & 0,94 \\ \hline
    Email         & 0,953 & 0,96 & 0,955 & 0,96 & 0,961 & 0,98 & 0,983 & 0,99 \\ \hline
    Form          & 0,718 & 0,76 & 0,804 & 0,85 & 0,848 & 0,88 & 0,902 & 0,94 \\ \hline
    Letter        & 0,671 & 0,71 & 0,805 & 0,83 & 0,835 & 0,86 & 0,898 & 0,92 \\ \hline
    Memo          & 0,758 & 0,79 & 0,844 & 0,90 & 0,879 & 0,90 & 0,920 & 0,93 \\ \hline
    News          & 0,709 & 0,67 & 0,890 & 0,89 & 0,883 & 0,90 & 0,938 & 0,93 \\ \hline
    Note          & 0,621 & 0,62 & 0,771 & 0,83 & 0,780 & 0,85 & 0,876 & 0,89 \\ \hline
    Report        & 0,495 & 0,43 & 0,527 & 0,61 & 0,638 & 0,71 & 0,735 & 0,81 \\ \hline
    Resume        & 0,969 & 0,97 & 0,784 & 0,80 & 0,803 & 0,96 & 0,951 & 0,97 \\ \hline
    Scientific    & 0,587 & 0,57 & 0,554 & 0,62 & 0,583 & 0,68 & 0,718 & 0,79 \\ \hline
    \textbf{Macro F1} & 0,702 & 0,71 & 0,784 & 0,82 & 0,811 & 0,86 & 0,885 & 0,91 \\ \hline
    \textbf{OA}       & 0,727 & 0,738 & 0,809 & 0,845 & 0,834 & 0,878 & 0,898 & 0,921 \\ \hline
  \end{tabular}
  \sumber{Eksekusi kelompok; kolom P dari \textcite{audebert2020}.}
\end{table}

Selain model di atas, studi pendahuluan juga menjalankan tiga uji lanjutan, semuanya
pada seed 42:
\begin{enumerate}
  \item fusion konkatenasi tanpa \textit{early stopping}: OA 0,8304, macro F1 0,8108;
  \item diagnostik fitur cabang teks, tanpa pelatihan;
  \item fusion yang diperbaiki: OA 0,8110, macro F1 0,8108.
\end{enumerate}
Ketiganya diringkas pada Subbab~\ref{sec:studi-pendahuluan}.
